\chapter{O sinal de alarme}
% versão completa: chapters/14_pain.tex

\pencilfig{14}{\pencilart{14}}{A dor é um alarme de incêndio, não um termômetro.}

Todos os outros sentidos comunicam como é o mundo. A dor, não. Ela comunica que algo ameaça o corpo, e o seu volume é, em grande parte, definido pela própria máquina, não pelo estímulo. A dor é um alarme de incêndio, não um termômetro.

Os sensores de dor só disparam com coisas fortes: calor acima de quarenta e poucos graus, pressão forte, substâncias corrosivas. Eles se acostumam bem menos que os sensores de toque --- o alarme não deve se calar enquanto a ameaça existe.

\section*{Dois fios}

Lembre de quando você bateu o dedo. Primeiro, uma dor aguda, precisa: onde e quando. Um instante depois, uma dor surda, difusa, demorada: quão ruim foi. São dois fios diferentes. Um é rápido, o seu sinal chega em dezenas de milissegundos. O outro é lento, o seu sinal leva cerca de um segundo.

\section*{A porta}

Por que, depois de uma pancada, esfregamos o lugar? Na medula espinhal, o sinal de dor passa por uma ``porta''. O toque acorda um neurônio inibitório que a fecha em parte. Por isso esfregar o lugar da pancada ajuda de verdade --- mas só contra dor moderada; a forte passa mesmo assim. Que a dor forte, nesse processo, desliga ela mesma o neurônio inibitório faz parte da teoria original da porta, e ainda não há confirmação direta disso.

\pencilfig{126}{%
% gate: the pain wire drives the relay neuron; touch wakes an inhibitory neuron that closes the gate
\draw[pen=pcAmber] (0,1.35) -- (3.05,1.0); \draw[arr=pcAmber] (2.8,1.03) -- (3.08,1.0);
\node[lbl, text=pcAmber!70!black, anchor=east] at (-0.05,1.35) {dor};
\draw[pen=pcBlue] (0,-0.15) -- (1.75,-0.15); \draw[arr=pcBlue] (1.5,-0.15) -- (1.78,-0.15);
\node[lbl, text=pcBlue, anchor=east] at (-0.05,-0.15) {toque};
\draw[dotpen=pcRed, wash=pcRed] (2.05,-0.15) circle (0.26);
\draw[pcRed, line width=0.7pt, -{Bar[width=6pt]}] (2.25,0.05) -- (3.2,0.66);
\node[lbl, text=pcRed, anchor=north] at (2.05,-0.45) {freio};
\draw[dotpen=pcGray, wash=pcGray] (3.4,0.9) circle (0.34);
\node[lbl, anchor=south] at (3.4,1.3) {porta};
\draw[pen] (3.75,0.9) -- (5.6,0.9); \draw[arr] (5.35,0.9) -- (5.65,0.9);
\node[lbl, anchor=west] at (5.7,0.9) {ao cérebro};
}{A porta na medula espinhal. O sinal de dor chega ao cérebro através de um neurônio-porta, e o toque acorda um neurônio inibitório que o fecha em parte. Por isso esfregar a pancada ajuda, mas só contra dor moderada.}

\section*{O botão de volume vindo de cima}

O cérebro sabe deixar a dor mais alta ou mais baixa por conta própria. De cima, do encéfalo, descem até a porta canais de rádio --- \nt{SERO}, \nt{NORA} e substâncias especiais, parecidas com analgésicos, que o próprio cérebro produz. Numa situação perigosa, a dor pode quase sumir: quem foge do perigo não tem tempo de mancar. A expectativa de alívio também abafa a dor, em parte pelas mesmas substâncias.

\section*{Quando o alarme emperra}

A dor repetida aumenta a sensibilidade: a mesma pancada passa a doer mais. Isso também é aprendizado --- um laço que se reforça sozinho enquanto a ferida cicatriza. Mas um laço forte tem um perigo: pode ``emperrar'' e continuar tocando quando a ferida já sarou --- como o grupo de neurônios do segundo capítulo, que se mantém sozinho.

\section*{Professor e interrupção}

A dor também é professora: para o ``melhor ou pior do que o esperado'' da dopamina, ela funciona como uma recompensa negativa. E ela interrompe o que está em curso, como uma rajada de noradrenalina: largue tudo, cuide do perigo. Já a reação mais rápida --- tirar a mão de algo quente --- se fecha direto na medula espinhal, com aquele mesmo par de músculos antagonistas do capítulo anterior, antes mesmo de você sentir a dor.

\fullversion{14}
